% Unit 068 terjemahan Bahasa Indonesia dari chapter5.tex baris 578--739.
% Sumber: Methods of Algebra, Volume 2, commit
% 9a5803ff77dd3257484cb177f851a73770a59dd3 (CC BY 4.0).
% stable-unit-id: o014.aljabr2.chapter5.filtered-complex-spectral-sequence
% Cakupan: Bagian 5.5 lengkap.

% segment-id: o014.li2.chapter5.filtered-complex-spectral-sequence.h001
\section{Barisan spektral kompleks terfiltrasi}\label{sec:filtered-cplx-ss}
\hypertarget{o014.aljabr2.chapter5.filtered-complex-spectral-sequence}{ }
\index{kompleks!terfiltrasi (filtered)}
\index{barisan spektral!dari kompleks terfiltrasi}

% segment-id: o014.li2.chapter5.filtered-complex-spectral-sequence.p001
Kita melanjutkan gagasan dari \S\ref{sec:filtered-diff-ss}, dengan kategori
abelian $\mathcal{A}$ tetap. Hasil
Proposisi~\ref{prop:filtered-diff-E} mengenai objek diferensial terfiltrasi
dapat diperluas agar morfisme $d$ boleh mempunyai derajat. Secara khusus,
teori ini dapat diterapkan pada kompleks terfiltrasi.

% segment-id: o014.li2.chapter5.filtered-complex-spectral-sequence.p002
Secara konkret, gantilah $\mathcal{A}$ sebelumnya dengan
$\mathcal{A}^{\Z}$ dan tuliskan $T$ untuk funktor translasinya. Menurut
Contoh~\ref{eg:cplx-as-diff}, objek $(X,d)$ dalam
$(\mathcal{A}^{\Z},T)_d$ dapat diidentifikasi dengan kompleks
$X\in\Obj(\cate{C}(\mathcal{A}))$. Jika selanjutnya kita mengambil filtrasi
menurun $(X,d,\mathrm{F}^\bullet X)$, maka
\begin{align*}
	\Hm(X,d)&=\left(\Hm^n(X)\right)_{n\in\Z}, \\
	\Hm\left(\mathrm{F}^pX,d\right)
	&=\left(\Hm^n\left(\mathrm{F}^pX\right)\right)_{n\in\Z}, \\
	\mathrm{F}^p\Hm^n(X)&:=
	\Image\left[\mathrm{F}^pX^n\cap\Ker(d_X^n)\to\Hm^n(X)\right]
	\quad\text{(Definisi~\ref{def:induced-filtration-H}).}
\end{align*}

% segment-id: o014.li2.chapter5.filtered-complex-spectral-sequence.p003
Jadi, $E_r$ dalam barisan spektral $\mathscr{E}$ yang dibangun sebelumnya
sebenarnya bernilai dalam $\mathcal{A}^{\Z\times\Z}$ dan merupakan objek
bergradasi ganda. Selain derajat filtrasi $p$, terdapat ``derajat internal''
$n$ yang berasal dari struktur kompleks; funktor translasi $S$ dan $T$ yang
bersesuaian dengan keduanya komutatif secara ketat. Kita telah mengetahui
bahwa derajat $d_r$ terhadap $S$ adalah $r$; sekarang kita tentukan derajatnya
terhadap $T$.

% segment-id: o014.li2.chapter5.filtered-complex-spectral-sequence.p004
\begin{itemize}
	\item Dalam pasangan eksak $\mathscr{C}_{(1)}$ yang dipakai untuk membangun
	$\mathscr{E}$, morfisme $\nwarrow$ berderajat $1$ terhadap $T$, sedangkan
	morfisme lainnya berderajat nol. Memang, $\nwarrow$ berasal dari morfisme
	penghubung kohomologi sehingga berderajat $1$ terhadap $T$, dan panah
	lainnya jelas berderajat nol.
	\item Dengan memakai deskripsi dalam Lema~\ref{prop:exact-couple-Z-B},
	pernyataan tersebut berlaku secara rekursif bagi $\mathscr{C}_{(r)}$ untuk
	semua $r$.
	\item Akibatnya, $d_r$ dan $d$ dalam $(X,d)$ sama-sama merupakan morfisme
	berderajat $1$ terhadap $T$. Hal ini juga langsung tampak dari deskripsi
	dalam Proposisi~\ref{prop:differential-exact-couple}.
\end{itemize}

% segment-id: o014.li2.chapter5.filtered-complex-spectral-sequence.p005
Untuk keperluan penerapan, secara konvensional kita mengganti indeks dengan $p$
dan $q:=n-p$. Dengan demikian,
\begin{gather*}
	E_r=(E_r^{p,q})_{(p,q)\in\Z^2}
	=\left(Z_r^{p,q}/B_r^{p,q}\right)_{(p,q)\in\Z^2}, \\
	E_0^{p,q}=\left(\gr^pX\right)^{p+q}, \\
	E_1^{p,q}=\Hm^{p+q}\left(\gr^pX,\gr^pd\right), \\
	d_r=\left(d_r^{p,q}:E_r^{p,q}\to E_r^{p+r,q-r+1}\right)_{(p,q)\in\Z^2}
	:E_r\xrightarrow{(r,-r+1)}E_r.
\end{gather*}
Dengan kata lain, kompleks terfiltrasi menghasilkan barisan spektral
bergradasi ganda kohomologis sebagaimana dalam
Definisi~\ref{def:graded-spectral-sequence}. Secara dual, kompleks rantai
dengan filtrasi menaik menghasilkan barisan spektral bergradasi ganda
homologis.

% segment-id: o014.li2.chapter5.filtered-complex-spectral-sequence.e001
\begin{example}\label{eg:canonically-bdd-ss}
	Perhatikan kompleks terfiltrasi $(X,d,\mathrm{F}^\bullet X)$. Jika untuk
	semua $n\in\Z$ berlaku $\mathrm{F}^0X^n=X^n$ dan
	$\mathrm{F}^{n+1}X^n=0$, maka $E_0$ terletak di kuadran pertama
	(Definisi~\ref{def:quadrant-ss}). Hal ini langsung mengikuti dari
	$E_0^{p,q}=(\gr^pX)^{p+q}$.
\end{example}

% segment-id: o014.li2.chapter5.filtered-complex-spectral-sequence.t001
\begin{proposition}\label{prop:filtered-cplx-ss}
	Untuk suatu kompleks $X\in\Obj(\cate{C}(\mathcal{A}))$ dengan filtrasi
	menurun $\mathrm{F}^\bullet X$, barisan spektral terkait memenuhi
	\begin{align*}
		Z_r^{p,q}&=
		\dfrac{\left(\mathrm{F}^pX^{p+q}\cap
		d^{-1}\mathrm{F}^{p+r}X^{p+q+1}\right)+\mathrm{F}^{p+1}X^{p+q}}
		{\mathrm{F}^{p+1}X^{p+q}}, \\
		B_r^{p,q}&=
		\dfrac{\left(\mathrm{F}^pX^{p+q}\cap
		d\mathrm{F}^{p-r+1}X^{p+q-1}\right)+\mathrm{F}^{p+1}X^{p+q}}
		{\mathrm{F}^{p+1}X^{p+q}},
	\end{align*}
	dan $d_r^{p,q}:E_r^{p,q}\to E_r^{p+r,q-r+1}$ diinduksi oleh
	\[ d^{-1}\mathrm{F}^{p+r}X^{p+q+1}\xrightarrow{d}
	\mathrm{F}^{p+r}X^{p+q+1}\cap\Ker(d). \]
\end{proposition}
\begin{proof}
	Dalam Proposisi~\ref{prop:filtered-diff-E}, sertakan derajat kompleks
	$n=p+q$ dan perhatikan bahwa $d$ berderajat $1$ terhadap funktor translasi
	kompleks $T$.
\end{proof}

% segment-id: o014.li2.chapter5.filtered-complex-spectral-sequence.p006
Dengan cara yang sama, \eqref{eqn:filtered-diff-infty} mempunyai versi
bergradasi ganda
\begin{equation}\label{eqn:filtered-cplx-infty}
	E_\infty^{p,q}=\dfrac{Z_\infty^{p,q}}{B_\infty^{p,q}}
	=\dfrac{\bigcap_r\left((\mathrm{F}^pX^{p+q}\cap
	d^{-1}\mathrm{F}^{p+r}X^{p+q+1})+\mathrm{F}^{p+1}X^{p+q}\right)}
	{\bigcup_r\left((\mathrm{F}^pX^{p+q}\cap
	d\mathrm{F}^{p-r+1}X^{p+q-1})+\mathrm{F}^{p+1}X^{p,q}\right)},
\end{equation}
dengan syarat $\bigcup$ dan $\bigcap$ yang ditampilkan ada. Dalam keadaan ini
$E_\infty=(E_\infty^{p,q})_{(p,q)\in\Z^2}$.

% segment-id: o014.li2.chapter5.filtered-complex-spectral-sequence.p007
Definisi--Proposisi~\ref{def:diff-obj-convergence} sekarang menjadi versi
bergradasi ganda: untuk semua $p,q$, objek $\gr^p\Hm^{p+q}(X)$ secara kanonik
dapat direalisasikan sebagai subkuosien dari $E_\infty^{p,q}$. Sifat
konvergensi barisan spektral pun diperinci sesuai dengan gradasi tersebut.

% segment-id: o014.li2.chapter5.filtered-complex-spectral-sequence.d001
\begin{definition}\label{def:filtered-cplx-convergence}
	\index{barisan spektral!konvergen lemah, konvergen kuat
	(weakly convergent, strongly convergent)}
	Untuk kompleks terfiltrasi $(X,d,\mathrm{F}^\bullet X)$ pada $\mathcal{A}$,
	bangun barisan spektral bergradasi kohomologis terkait $(E_r,d_r)_r$. Jika
	$\bigcap_r$ dan $\bigcup_r$ dalam \eqref{eqn:filtered-cplx-infty} ada, maka
	$\gr^p\Hm^{p+q}(X)$ secara kanonik dapat direalisasikan sebagai subkuosien
	dari $E_\infty^{p,q}$.
	\begin{itemize}
		\item Jika $\gr^p\Hm^{p+q}(X)=E_\infty^{p,q}$ untuk semua
		$(p,q)\in\Z^2$, maka $(E_r,d_r)_r$ disebut \emph{konvergen lemah}.
		\item Jika barisan tersebut konvergen lemah dan, untuk semua $n\in\Z$,
		filtrasi $\mathrm{F}^\bullet\Hm^n(X)$ menyeluruh serta lengkap, maka
		$(E_r,d_r)_r$ disebut \emph{konvergen kuat}.
	\end{itemize}
\end{definition}

% segment-id: o014.li2.chapter5.filtered-complex-spectral-sequence.c001
\begin{convention}\label{con:ss-conv}
	\index[sym1]{Erpq abuts to H@$E_r^{p,q}\Rightarrow\Hm^{p+q}(X)$}
	Konvergensi kuat barisan spektral juga ditulis
	$E_r^{p,q}\Rightarrow\Hm^{p+q}(X)$. Subskrip $r$ biasanya ditulis secara
	konkret sebagai $1$, $2$, dan seterusnya, bergantung pada halaman barisan
	spektral yang hendak kita deskripsikan.

	Lebih umum, misalkan diberikan barisan spektral bergradasi ganda
	kohomologis $\mathscr{E}$, objek bergradasi terfiltrasi
	$(H,\mathrm{F}^\bullet H)$, dan isomorfisme
	$E_\infty^{p,q}\simeq\gr^pH^{p+q}$, dengan
	$\mathrm{F}^\bullet H$ menyeluruh serta lengkap. Situasi ini juga kita
	nyatakan dengan $E_r^{p,q}\Rightarrow H^{p+q}$. Kasus barisan spektral
	bergradasi ganda homologis diperlakukan dengan cara serupa; khususnya,
	$E^r_{p,q}\Rightarrow H_{p+q}$ berarti
	$E^\infty_{p,q}\simeq\gr_pH_{p+q}$.
\end{convention}

% segment-id: o014.li2.chapter5.filtered-complex-spectral-sequence.p008
Teorema konvergensi klasik berikut sudah cukup untuk penerapan awal. Untuk
syarat konvergensi yang lebih umum, lihat \cite{Boa99}.

% segment-id: o014.li2.chapter5.filtered-complex-spectral-sequence.t002
\begin{theorem}[Teorema konvergensi klasik]\label{prop:classical-convergence}
	Perhatikan kompleks terfiltrasi $(X,d,\mathrm{F}^\bullet X)$. Andaikan
	bahwa untuk setiap $n\in\Z$,
	\begin{compactitem}
		\item filtrasi $\mathrm{F}^\bullet X^n$ menyeluruh
		(Definisi~\ref{def:filtration-properties}),
		\item terdapat $N=N(n)$ dengan $\mathrm{F}^NX^n=0$.
	\end{compactitem}
	Maka barisan spektral terkait konvergen kuat
	(Definisi--Proposisi~\ref{def:diff-obj-convergence}).

	Jika syarat tersebut diperkuat sehingga filtrasi setiap $X^n$ hingga
	(Definisi~\ref{def:filtration}), maka filtrasi terinduksi pada $\Hm^n(X)$
	juga hingga. Dalam keadaan ini barisan spektral terkait terbatas
	(Definisi~\ref{def:bdd-ss}).
\end{theorem}
\begin{proof}
	Menurut versi bergradasi Lema~\ref{prop:induced-filtration-H}, filtrasi
	terinduksi $\mathrm{F}^\bullet\Hm^n(X)$ menyeluruh serta lengkap. Jadi cukup
	dibuktikan konvergensi lemah.

	Kita perlu mengingat bukti
	Definisi--Proposisi~\ref{def:diff-obj-convergence}, yang menjelaskan cara
	merealisasikan $\gr^p\Hm^{p+q}(X)$ sebagai subkuosien dari
	$E_\infty^{p,q}$ untuk setiap $(p,q)\in\Z^2$. Unsur kuncinya adalah
	\begin{equation*}\begin{split}
		\bigcap_r&\left(\left(\mathrm{F}^pX^{p+q}\cap
		d^{-1}\mathrm{F}^{p+r}X^{p+q+1}\right)+\mathrm{F}^{p+1}X^{p+q}\right)
		\\
		&\supset(\mathrm{F}^pX^{p+q}\cap\Ker(d))+\mathrm{F}^{p+1}X^{p+q},\\
		\bigcup_r&\left(\left(\mathrm{F}^pX^{p+q}\cap
		d\mathrm{F}^{p-r+1}X^{p+q-1}\right)+\mathrm{F}^{p+1}X^{p+q}\right)
		\\
		&\subset(\mathrm{F}^pX^{p+q}\cap\Image(d))+\mathrm{F}^{p+1}X^{p+q}.
	\end{split}\end{equation*}
	Membuktikan konvergensi lemah setara dengan memperkuat kedua inklusi ini
	menjadi kesamaan. Akan tetapi, jika $r\gg0$ relatif terhadap $p,q$, isi
	$\bigcap_r$ pada baris pertama tepat merupakan
	$(\mathrm{F}^pX^{p+q}\cap\Ker(d))+\mathrm{F}^{p+1}X^{p+q}$, sehingga
	kesamaan berlaku.

	Untuk baris kedua, suku $+\mathrm{F}^{p+1}X^{p+q}$ dapat dikeluarkan dari
	$\bigcup_r$. Ingat bahwa $d(\mathrm{F}^\bullet X^n)$ merupakan filtrasi
	menyeluruh dari $d(X^n)$. Karena itu,
	\begin{multline*}
		\bigcup_r\left(\mathrm{F}^pX^{p+q}\cap
		d\mathrm{F}^{p-r+1}X^{p+q-1}\right)\\
		\xlongequal{\because\;\text{\eqref{eqn:exhaustion-intersection}}}
		\mathrm{F}^pX^{p+q}\cap\bigcup_r
		d\mathrm{F}^{p-r+1}X^{p+q-1}\\
		=\mathrm{F}^pX^{p+q}\cap d\bigcup_r
		\mathrm{F}^{p-r+1}X^{p+q-1}
		=\mathrm{F}^pX^{p+q}\cap\Image(d).
	\end{multline*}

	Terakhir, andaikan filtrasi setiap $X^n$ hingga. Maka
	$\mathrm{F}^\bullet\Hm^n(X)$ tentu terbatas. Untuk membuktikan bahwa
	barisan spektralnya terbatas, misalkan $r,n$ tetap dan ambil $p+q=n$. Menurut
	deskripsi dalam Proposisi~\ref{prop:filtered-cplx-ss}, jika $p\gg0$ maka
	pembilang $Z_r^{p,q}$ bernilai $0$, sedangkan jika $p\ll0$ maka
	penyebutnya sama dengan $X^n$. Jadi hanya ada berhingga banyak
	$(p,q)$ dengan $E_{r+1}^{p,q}\neq0$.
\end{proof}

% segment-id: o014.li2.chapter5.filtered-complex-spectral-sequence.p009
Versi dualnya jelas dan tidak kita ulangi.

% segment-id: o014.li2.chapter5.filtered-complex-spectral-sequence.t003
\begin{corollary}[Barisan eksak untuk suku berderajat rendah]
\label{prop:low-degree-ss}
	Misalkan kompleks terfiltrasi $(X,d,\mathrm{F}^\bullet X)$ memenuhi
	\[ \mathrm{F}^0X^n=X^n,\qquad \mathrm{F}^{n+1}X^n=0,
	\qquad n\in\Z, \]
	seperti dalam Contoh~\ref{eg:canonically-bdd-ss}. Barisan spektral terkait
	memberikan barisan eksak kanonik
	\[ 0\to E_2^{1,0}\to\Hm^1(X)\to E_2^{0,1}\xrightarrow{d}E_2^{2,0}
	\to\Hm^2(X). \]
	Secara dual, untuk kompleks rantai $X$ dengan filtrasi menaik, jika
	$\mathrm{F}_{-1}X=0$ dan $\mathrm{F}_nX=X$, terdapat barisan eksak kanonik
	\[ \Hm_2(X)\to E^2_{2,0}\xrightarrow{d}E^2_{0,1}\to\Hm_1(X)
	\to E^2_{1,0}\to0. \]
\end{corollary}
\begin{proof}
	Substitusikan $p=2$ dalam \eqref{eqn:low-exact-coh} untuk memperoleh
	barisan eksak
	\[ 0\to E_\infty^{0,1}\to E_2^{0,1}\xrightarrow{d}E_2^{2,0}
	\to E_\infty^{2,0}\to0. \]
	Berdasarkan Teorema Konvergensi Klasik~\ref{prop:classical-convergence},
	berlaku $E_\infty^{0,1}\simeq\gr^0\Hm^1(X)$,
	$E_\infty^{1,0}\simeq\gr^1\Hm^1(X)$, dan
	$E_\infty^{2,0}\simeq\gr^2\Hm^2(X)$. Dari syarat diperoleh
	\begin{gather*}
		\gr^2\Hm^2(X)=\mathrm{F}^2\Hm^2(X),\qquad
		\gr^1\Hm^1(X)=\mathrm{F}^1\Hm^1(X),\\
		\gr^0\Hm^1(X)=\Hm^1(X)/\mathrm{F}^1\Hm^1(X)
		=\Hm^1(X)/E_\infty^{1,0}.
	\end{gather*}
	Selain itu, substitusi $p=1$ dalam \eqref{eqn:edge-ss-coh} memberikan
	$E_2^{1,0}=E_\infty^{1,0}$. Menggabungkan semua kesamaan ini menghasilkan
	barisan eksak yang dinyatakan. Versi homologis serupa.
\end{proof}

% segment-id: o014.li2.chapter5.filtered-complex-spectral-sequence.p010
Jika $E_r^{p,q}$ konvergen kuat, maka setelah mengetahui cukup banyak halaman
$(E_r,d_r)$, pada prinsipnya kita dapat membaca
$(\gr^p\Hm^{p+q}(X))_{p,q\in\Z}$ dari $E_\infty$. Akan tetapi, beralih dari
$(\gr^p\Hm^n(X))_p$ ke $\Hm^n(X)$ berarti menentukan serangkaian ekstensi
dalam $\mathcal{A}$, yang umumnya tidak mudah kecuali $\Hm^n(X)$ terbelah.
Jika kita hanya mempertimbangkan sifat yang lebih kasar, barisan spektral
kadang-kadang memberi jawaban ringkas. Contoh berikut pada dasarnya merupakan
penerapan prinsip Euler--Poincar\'e.

% segment-id: o014.li2.chapter5.filtered-complex-spectral-sequence.l001
\begin{lemma}\label{prop:finite-degeneration}
	Misalkan $\mathscr{E}=(E_r,d_r)_{r\geq1}$ adalah barisan spektral
	bergradasi ganda kohomologis dan terdapat $r$ sedemikian sehingga
	\[ \left\{(p,q)\in\Z^2:E_r^{p,q}\neq0\right\}
	\quad\text{merupakan himpunan hingga}. \]
	Maka untuk $r\gg0$, $\mathscr{E}$ mengalami degenerasi pada halaman $E_r$.
\end{lemma}
\begin{proof}
	Menurut syarat tersebut, jika $r$ cukup besar maka semua suku taknol $E_r$
	terkonsentrasi pada suatu daerah dengan panjang dan lebar hingga. Terapkan
	Contoh~\ref{eg:SS-finite-width}.
\end{proof}

% segment-id: o014.li2.chapter5.filtered-complex-spectral-sequence.t004
\begin{proposition}\label{prop:EP-ss}
	Misalkan barisan spektral bergradasi ganda kohomologis $(E_r,d_r)_r$
	memenuhi syarat-syarat berikut:
	\begin{compactitem}
		\item terdapat $r$ sedemikian sehingga
		$\{(p,q)\in\Z^2:E_r^{p,q}\neq0\}$ merupakan himpunan hingga;
		\item berlaku konvergensi kuat $E_r^{p,q}\Rightarrow H^{p+q}$ dalam
		arti Konvensi~\ref{con:ss-conv}.
	\end{compactitem}
	Maka dalam grup $\mathrm{K}_0(\mathcal{A})$ dari Definisi~\ref{def:K0},
	untuk $r\gg0$ berlaku
	\[ \sum_{n\in\Z}(-1)^n[H^n]
	=\sum_{p,q\in\Z}(-1)^{p+q}[E_r^{p,q}], \]
	dan kedua ruas merupakan jumlah hingga.
\end{proposition}
\begin{proof}
	Untuk $r\gg0$, hanya berhingga banyak suku di ruas kanan yang taknol. Selain
	itu,
	\[ E_{r+1}^{p,q}\simeq\Hm\left[
	E_r^{p-r,q+r-1}\xrightarrow{d}E_r^{p,q}
	\xrightarrow{d}E_r^{p+r,q-r+1}\right]. \]
	Jika derajat dihitung menurut $n:=p+q$, maka $d$ berderajat $1$. Ambil
	jumlah berselang-seling terhadap $n$ dalam $\mathrm{K}_0(\mathcal{A})$,
	lalu terapkan Teorema~\ref{prop:EP}, untuk memperoleh
	\[ \sum_{p,q\in\Z}(-1)^{p+q}[E_r^{p,q}]
	=\sum_{p,q\in\Z}(-1)^{p+q}[E_{r+1}^{p,q}]. \]

	Lema~\ref{prop:finite-degeneration} menunjukkan bahwa barisan spektral pada
	akhirnya mengalami degenerasi. Jadi, jika $r\gg0$, berlaku
	$E_r^{p,q}=E_\infty^{p,q}$ untuk semua $(p,q)$. Karena itu,
	\[ \sum_{p,q}(-1)^{p+q}[E_r^{p,q}]
	=\sum_{p,q}(-1)^{p+q}[\gr^p\Hm^{p+q}(X)]
	=\sum_n(-1)^n\sum_p[\gr^pH^n], \]
	dan syarat-syarat di atas memastikan bahwa semua jumlah tersebut hingga.
	Menurut Lema~\ref{prop:K0-prep}~(v) dan sifat $\mathrm{F}^\bullet H$, kita
	juga mempunyai
	$\sum_p[\gr^pH^n(X)]=[\Hm^n(X)]$.
\end{proof}

% segment-id: o014.li2.chapter5.filtered-complex-spectral-sequence.p011
Jika $(E_r,d_r)_r$ berasal dari kompleks terfiltrasi
$(X,d,\mathrm{F}^\bullet)$ dan filtrasi $\mathrm{F}^\bullet X^n$ pada setiap
$X^n$ hingga, maka Teorema~\ref{prop:classical-convergence} memastikan bahwa
syarat konvergensi kuat $E_r^{p,q}\Rightarrow\Hm^{p+q}(X)$ dalam
Proposisi~\ref{prop:EP-ss} berlaku secara otomatis. Akan tetapi, syarat
mengenai $\{(p,q):E_r^{p,q}\neq0\}$ tidak otomatis berlaku.

% segment-id: o014.li2.chapter5.filtered-complex-spectral-sequence.p012
Kompleks terfiltrasi sudah cukup untuk menghasilkan beberapa barisan spektral
sederhana namun berguna; lihat latihan bab ini. Karena pokok bahasan buku ini
adalah aljabar, beberapa barisan spektral yang paling dikenal berasal dari
bikompleks. Karena itu, sekarang kita beralih ke kasus bikompleks.
